\subsection{DEFINITION OF COMPLEX NUMBERS}

The polynomial \(X^{2} + 1\) has no root in R, because \(x^{2} + 1 \geqslant 1\) for all \(x \in R\) ; it is therefore irreducible over R. {[}This is a particular case of the analogous result which applies to any ordered field.{]}

DEFINITION 1. The field \(\mathbf{R}[\mathbf{X}] / (\mathbf{X}^2 + 1)\) is called the field of complex numbers and is denoted by \(\mathbf{C}\) . The canonical image of \(\mathbf{X}\) in \(\mathbf{C}\) is denoted by \(i\) , so that \(\mathbf{C}\) is obtained from the field \(\mathbf{R}\) by algebraic adjunction of the root \(i\) of the polynomial \(\mathbf{X}^2 + 1\) . The elements of \(\mathbf{C}\) are called complex numbers.

From an algebraic point of view the importance of the field \(\mathbf{C}\) is due to the following fundamental theorem:

THEOREM 1. (d'Alembert-Gauss). The field \(\mathbf{C}\) of complex numbers is algebraically closed.

For the proof it is enough to establish that (i) every element \(\geqslant0\) in R has a square root, and \((\mathrm{ii})\) every polynomial of odd degree with coefficients in R has at least one root in R. The first of these assertions has already been proved in Chapter IV, § 3, no. 3. As to the second, if \(f(\mathbf{X}) = a_{0}\mathbf{X}^{n} + a_{1}\mathbf{X}^{n-1} + \cdots + a_{n}\) is a polynomial of odd degree \(n(a_{0} \neq 0)\) with real coefficients, we may write \(f(x) = a_{0}x^{n}g(x)\) for \(x \neq 0\) , where

\[
g (x) = \mathrm{I} + \frac {a _ {1}}{a _ {0} x} + \dots + \frac {a _ {n}}{a _ {0} x ^ {n}}
\]

tends to \(+1\) as \(x\) tends to \(+\infty\) or \(-\infty\) . Hence there is a number \(a > 0\) such that \(f(a)\) has the sign of \(a_0\) and \(f(-a)\) the sign of \(-a_0\) ;
and so by Bolzano's theorem (Chapter IV, § 6, no. 1, Theorem 2), \(f\) has at least one root in \([-a, a]\) .

Remarks. 1) Theorem 1 can be proved without invoking the theory of ordered fields by using properties of the topology of the field C, which will be defined below (no. 2); see § 2, Exercise 2 and also the part of this series devoted to algebraic topology, where the theorem of d'Alembert-Gauss will appear as a consequence of results on the degree of a mapping.

\begin{enumerate}
\def\labelenumi{\arabic{enumi})}
\setcounter{enumi}{1}
\tightlist
\item
  Since C is of degree 2 over R, it follows that C is, up to isomorphism, the only algebraic extension of R other than R itself, and that there is no field contained in C which contains R, other than R and C.
\end{enumerate}

We know that \(\mathbf{R}\) may be identified with a subfield of \(\mathbf{C}\) , and that every \(z \in \mathbf{C}\) can be written uniquely in the form \(x + iy\) , where \(x\) and \(y\) are real; \(x\) is called the real part of \(z\) and is denoted by \(\Re(z)\) ; \(y\) the imaginary part of \(z\) , denoted by \(\Im(z)\) . Complex numbers of the form \(iy\) ( \(y\) real) are called pure imaginary. The relation \(x + iy = o\) ( \(x, y\) real) is equivalent to \(x = o\) and \(y = o\) .

Since \(i^2 = -1\) , the elements of \(\mathbf{C}\) (when given by their real and imaginary parts) satisfy the following rules of calculation:

\begin{enumerate}
\def\labelenumi{(\arabic{enumi})}
\tightlist
\item
\end{enumerate}

\[
(x + i y) + (x ^ {\prime} + i y ^ {\prime}) = (x + x ^ {\prime}) + i (y + y ^ {\prime})\tag{2}
\]

\[
(x + i y) (x ^ {\prime} + i y ^ {\prime}) = (x x ^ {\prime} - y y ^ {\prime}) + i (x y ^ {\prime} + x ^ {\prime} y).
\]

In particular, \((x + iy)(x - iy) = x^{2} + y^{2}\in \mathbf{R}\) , so that, if \(x + iy\neq 0\) ,

\[
\frac {1}{x + i y} = \frac {x}{x ^ {2} + y ^ {2}} - i \frac {y}{x ^ {2} + y ^ {2}}.\tag{3}
\]

The second root of the polynomial \(X^{2}+1\) in C is -i; consequently the only automorphism of C, other than the identity mapping, which leaves all real numbers invariant, is that which maps \(z=x+iy\) to x-iy; the latter is denoted by \(\bar{z}\) and (in agreement with the general definitions) is called the complex number conjugate to z. We have

\[
\mathcal {R} (z) = \frac {\mathrm{I}}{2} (z + \bar {z}) \quad \text { and } \quad \mathfrak {J} (z) = \frac {\mathrm{I}}{2 i} (z - \bar {z}).
\]

By reason of this automorphism, if \(f(z)\) is a polynomial with real coefficients, we have \(f(\overline{z}) = \overline{f(z)}\) for all \(z \in \mathbf{C}\) .

The real number \(z\bar{z} = x^2 + y^2\) is called the algebraic norm of \(z\) , or simply the norm whenever there is no risk of confusion; it is a real number \(\geqslant 0\) , which vanishes only if \(z = 0\) . The real number \(\sqrt{z\bar{z}} = \sqrt{x^2 + y^2} \geqslant 0\) reduces to the absolute value of \(z\) when \(z\) is real, and we still call it the absolute value of \(z\) , and denote it by \(|z|\) , when \(z\) is any complex number. The relation \(|z| = 0\) is equivalent to \(z = 0\) . If \(z\) and \(z'\) are two complex numbers, the conjugate of \(zz'\) is \(\overline{zz'}\) , hence \(|zz'|^2 = zz' \overline{zz'} = |z|^2 |z'|^2\) and therefore \(|zz'| = |z|. |z'|\) : the absolute value of a product is the product of the absolute values of the factors. In particular, if \(z \neq 0\) and \(z' = 1/z\) , we have \(|1/z| = 1/|z|\) .

Finally, for all complex numbers \(z, z'\) , we have the triangle inequality

\[
\left| z + z ^ {\prime} \right| \leqslant | z | + \left| z ^ {\prime} \right|.\tag{4}
\]

